\chapter{Introduction}
\label{ch:introduction}

\section{Background and Motivation}
\label{sec:background_motivation}

\glspl{dt} have become a standard component in modern industrial automation for connecting physical assets with their virtual counterparts. By maintaining accurate virtual copies of systems like industrial robots, \glspl{dt} enable real-time monitoring, predictive maintenance, and remote control. In time-critical \gls{cps} applications, the utility of a \gls{dt} is strictly defined by the synchronization accuracy between the physical robot and the digital model~\cite{9013428}. Consequently, ensuring state consistency typically demands \gls{urllc} capabilities to support high-frequency updates.

However, achieving this synchronization through standard periodic transmission is often inefficient in bandwidth-constrained environments. In many control architectures, the \gls{dt} acts as a global supervisor that maintains a high-fidelity state estimate. When the \gls{dt} transmits full state updates to the robot at every time step, it results in significant data redundancy. Since the robot already possesses local sensors to maintain a short-term estimate, "blindly" broadcasting the entire state vector from the \gls{dt} consumes scarce network resources without providing new actionable information. This mismatch between the \gls{dt}'s global knowledge and the robot's local perception creates an opportunity to replace continuous payload reconstruction with sparse, verification-based communication.

To mitigate this redundancy, recent research proposes identification-based \gls{id} synchronization~\cite{9789759}. Unlike classical Shannon-type communication, which aims to reconstruct the full message, \gls{id} schemes focus on a verification task: confirming whether the remote state matches a local hypothesis. Building on this concept, prior works~\cite{fi16030078} demonstrate that transmitting a compressed ``verifier'' (e.g., a short hash or randomized code) instead of the full state can significantly reduce communication traffic. This aligns with semantic communication in \gls{sixg} networks~\cite{doi:https://doi.org/10.1002/9781119765554.ch16}, shifting the focus from what data is sent to what value the data holds for control.

In real-time control loops, deploying \gls{id} synchronization exposes two engineering trade-offs that motivate this thesis. First, the computational overhead of encoding on edge devices may offset the time saved by reducing transmission traffic. This shift from a communication-bound to a computation-bound bottleneck calls for an end-to-end latency analysis that accounts for both compute and communication. Second, reliability becomes a primary concern when moving from exact to semantic verification. In practice, enforcing bitwise synchronization (standard \gls{id}) can be overly conservative, triggering unnecessary updates for benign deviations such as sensor noise. This motivates semantic verification over ``safety sets'', i.e., \gls{kid}. However, the relaxation introduces collision risks, specifically Type~II (miss) errors, where unsafe desynchronizations are incorrectly classified as admissible.

\section{Problem Statement}
\label{sec:problem_statement}

Existing research on \gls{id} has primarily focused on information-theoretic capacity or simple traffic suppression, lacking a comprehensive evaluation of latency, physical safety, and computational feasibility. Specifically, it remains unclear under what bandwidth and computational constraints \gls{id} truly outperforms traditional transmission, and how to quantify the Pareto trade-off between semantic tolerance and safety risk in \gls{kid} systems.

Moreover, semantic tolerance introduces structural constraints that are not visible in purely static formulations: as the acceptance set grows with tolerance, the finite verifier space $2^{n_{ver}}$ can become too small to distinguish all tolerated states, imposing a feasibility limit on the achievable miss probability. Characterizing these feasibility limits and their implications in dynamic settings is part of the open problem addressed in this thesis.

\section{Contributions}
\label{sec:contributions}

To address these limitations, this thesis proposes a comprehensive framework ranging from theoretical modeling to dynamic evaluation, investigating not only the static identification problem but also error evolution in dynamic environments.

First, we develop an end-to-end latency model combining computational and transmission delays to compare algebraic coding (\gls{rsid}) and truncated hashing based on \glsentryshort{sha}-256. For the evaluated software stack, truncated \glsentryshort{sha}-256 achieves low and largely payload-independent encoding latency, while \gls{rsid} incurs higher computation that depends on payload symbolization. These results quantify a practical compute--communication trade-off for edge deployments, while also motivating \gls{rsid} for safety-critical deployments requiring rigorous mathematical guarantees (e.g., provable upper bounds under stated assumptions) on worst-case error-detection failure probability.

Second, we generalize identification to \gls{kid}, introducing semantic tolerance to filter benign jitter. Through parameter sweeps, we map the risk--speedup Pareto frontier and show that moderate tolerance can substantially reduce fallback frequency while keeping safety risks bounded. The analysis also indicates a knee point where additional tolerance yields diminishing speedup while reducing verifier selectivity.

Third, the evaluation is extended to a time-correlated random-walk drift model to capture risks that static analysis may overlook. The results demonstrate that verifiers with insufficient bit-depth (e.g., $n_{ver}=8$) can transform isolated collisions into persistent, undetected boundary breaches once the state drifts beyond the tolerance $R$. To quantify synchronization throughput, we analyze full-state recovery timing via the \gls{sri}. The dynamic results separate the roles of the parameters: $R$ mainly determines how often boundary crossings trigger recovery attempts (updates become sparser as $R$ increases), while $n_{ver}$ determines how often a triggered recovery is missed due to collisions.

\section{Thesis Organization}
\label{sec:organization}

The remainder of this thesis is organized as follows. Chapter~\ref{ch:background} reviews related work, covering the transition from data reconstruction to semantic identification, the mathematical principles of encoding mechanisms, and the control-theoretic background. Chapter~\ref{ch:design_methodology} details the system model, data flow design, and the derivation of theoretical error bounds for \gls{kid}. Chapter~\ref{sec:implementation} describes the implementation of the modular Python simulation framework and the experimental design strategy. Chapter~\ref{sec:results_and_evaluation} presents the experimental results, including latency sensitivity analysis, Pareto optimization of \gls{kid}, and risk assessment under dynamic drift scenarios. Finally, Chapter~\ref{sec:conclusion_and_future_work} concludes the thesis and outlines future research directions.